\documentclass{article}
\usepackage{fontspec}
\usepackage{xeCJK}
\usepackage{graphicx}
\usepackage{hyperref}
\usepackage{amsmath}
\usepackage{amssymb}
\usepackage{ulem}
\usepackage{geometry}
\usepackage{color}
\usepackage{fancyhdr}
\usepackage{lastpage}
\usepackage{booktabs}
\usepackage{caption}

% 英文正文字体
\setmainfont{Times New Roman}
\setsansfont{Arial}
\setmonofont{Consolas}

% 中文字体
\setCJKmainfont{SimSun}
\setCJKsansfont{Microsoft YaHei}
\setCJKmonofont{FangSong}

\title{未命名文档}
\author{}
\date{\today}

\begin{document}
\maketitle

\section{引言}
请在此处撰写引言内容。

\section{主体}
请在此处撰写主体内容。

\section{格式示例与交叉引用}
本小节展示行内公式、带编号公式、图表、文献引用与脚注的用法，可根据需要修改或删除。

\subsection{公式}
行内公式示例：爱因斯坦质能关系 $E=mc^2$，勾股定理 \( a^2 + b^2 = c^2 \)。带编号的块级公式示例：
\begin{equation}
  \label{eq:example}
  \int_{-\infty}^{+\infty} e^{-x^2} \, dx = \sqrt{\pi}.
\end{equation}
文中可引用公式“见式~\eqref{eq:example}”或“式~\ref{eq:example}”。

\subsection{插图}
将图片放在 \texttt{figures} 目录下，用 \texttt{includegraphics} 插入。图~\ref{fig:example} 为占位示意。
\begin{figure}[htbp]
  \centering
  \rule{4cm}{2.5cm}  % 占位，替换为 \includegraphics[width=4cm]{figures/example.pdf}
  \caption{示例图标题。}
  \label{fig:example}
\end{figure}

\subsection{表格}
表~\ref{tab:example} 为表格示例，可用 \texttt{booktabs} 绘制三线表。
\begin{table}[htbp]
  \centering
  \caption{示例表标题。}
  \label{tab:example}
  \begin{tabular}{lcc}
    \toprule
    项目 & 数值一 & 数值二 \\
    \midrule
    A    & 1.0   & 2.0   \\
    B    & 3.0   & 4.0   \\
    \bottomrule
  \end{tabular}
\end{table}

\subsection{文献引用与脚注}
参考文献在文末 \texttt{thebibliography} 中列出，文中用 \cite{ref1}、\cite{ref2} 引用。脚注示例：这是一段正文\footnote{这是脚注内容，用于补充说明。}，可有多处脚注。

\section{结论}
请在此处撰写结论内容。

% 参考文献（可按需增删 \bibitem）
\begin{thebibliography}{99}
\bibitem{ref1} 作者. 题名[J]. 刊名, 年, 卷(期): 页码.
\bibitem{ref2} 作者. 书名[M]. 出版地: 出版社, 年.
\end{thebibliography}

\end{document}
